\documentclass[12pt]{article}
\usepackage[margin=1in]{geometry}
\usepackage{tgtermes}
\usepackage{setspace}
\onehalfspacing
\usepackage{microtype}
\usepackage{amsmath,amssymb}
\usepackage[hidelinks,pdftitle={Robust solvability is not a realized procedure},pdfauthor={}]{hyperref}
\usepackage[authoryear,round]{natbib}
\setcitestyle{aysep={}}
\pagestyle{plain}
\begin{document}
\begin{center}
{\Large Robust solvability is not a realized procedure\par}
\vspace{0.75em}
\end{center}

\begin{abstract}
A finite wiring can be robust in a precise procedural sense: the test that asks whether a tuple of local maps has a unique history returns the same kind of answer, namely exactly one history, for every tuple. A further theorem says that the same condition is necessary and sufficient for a canonical closed loop to stay unitary under every local unitary and every matched finite ancilla. That is still a fact about a procedure defined on the wiring. It is not the convergence of independent detections of a phenomenon, and it is not evidence that a physical system implements the procedure. The second gap is the one Karaca marks between the robustness of a procedure and the robustness of a system that would have to carry it out. A bare reversible relabeling exists for wirings that fail the test. A finite search can decide the test for an explicit table and is a different procedure. A known distribution, read in one basis from a known board, falls outside the convex hull of three-bit process functions if an external enumeration is granted; that calculation is not an experiment. The mathematics is an alternative adaptation of results already public. It is not offered as an independent study, and it does not transfer the warrant of one robustness notion to another.
\end{abstract}

\noindent\textbf{AI statement.} The writing of this paper involved generative assistance. A full account is in Appendix~\ref{app:ai}. The account does not meet every reconstructability demand in the journal's current FAQ: the tool version is unrecorded, and a complete verbatim log of intermediate prompts is not retained.

\section{The question}
\label{sec:q}

The finite results used below are already public. They say when a wiring of finite alphabets has exactly one history for every replacement of its local maps, and when a particular unitary extension of that wiring stays unitary if every party inserts a unitary on a matched finite ancilla. They also separate, in one computational basis, a known distribution from the convex hull of three-bit process functions, subject to checks that are not all written out again here. Nothing in this paper adds a theorem.

The philosophical question is what kind of scientific virtue those results are. They are easy to praise under the word ``robustness,'' and easy, once praised, to treat as an explanation of a target process or as a reason to think the process has been implemented. The praise is cheap because several different virtues share the word. \citet{kuorikoski2007}, reporting \citet{woodward2006}, distinguish inferential robustness, derivational robustness, measurement robustness, and causal robustness, and they record Woodward's warning that arguments which support one notion do not automatically support another. I have not re-read Woodward's article. I use the warning only as they state it, and I do not attribute to him any further typology.

The notion that fits the theorems is closer to what \citet{karaca2022} calls procedure robustness than to the convergence of independent results. Karaca defines procedure robustness as ``the capacity of an experimental procedure to maintain its intended function invariant despite possible variations in its inputs.'' Result robustness, in his terms, is the existence of convergent experimental results obtained through different and independent means of detection. He argues that procedure robustness is a precondition for result robustness and is not sufficient for it. He also separates both from what he calls system robustness, the capacity of a system to maintain its functioning under perturbation. A procedure can be robust, in his sense, and still fail to be carried out, because the system that would have to implement it is not adequate. That last distinction is the one this paper needs. It is drawn from his text. It is not a slogan added after the fact.

The claim, then, is narrow. Universal consistency is procedure robustness of a solvability test. The canonical unitary equivalence is the same robustness claim for a larger class of inputs, the local unitaries. Neither claim is result robustness, neither is an inductive leap from the model to a target, and neither is implementation. The paper is an alternative adaptation of one mathematical note to this question. A reader who wants the proofs should use the note. A reader who wants a different philosophical question---modularity as against directed causal order, or admissibility as against model explanation in Bokulich's sense---is reading a different adaptation, not a different result.

\section{Three notions that do not share a warrant}
\label{sec:three}

Karaca's article is about experiments in high-energy physics, not about process functions. The transfer has to be argued, not assumed. What I take from him is the shape of the definitions, together with his explicit refusal to collapse procedure into system.

On his account, result robustness is invariance of an outcome across independent means of detection. He illustrates it with the \(K^{+}_{e2}\) branching-ratio measurement, where two methods of analysis gave compatible ratios, and with the CDF top-quark analyses, where agreement among tagging algorithms was used, on Staley's telling as Karaca reports it, as second-order evidence that a procedure was doing what it was supposed to do, rather than as direct evidence for the quark. I have not re-read Staley or Franklin. The cases matter here only as Karaca uses them: convergence of results is one epistemic achievement, and support for the reliability of a procedure is another. He then defines procedure robustness by a different invariance. The robust aspect is not the result. It is the procedure's maintenance of its intended function under variation of inputs. The function is whatever the procedure was designed to do. The variations are variations in what is fed to it.

He is also explicit that this is not yet a claim about the machine that runs the procedure. In the technical literature he surveys, robustness is often a system property: the system keeps working despite fluctuation in its parts or its environment. Karaca says that if an experimental system lacks that property, ``this does not necessarily mean that the procedure to be implemented by this experimental system would lack sufficient'' procedure robustness. ``Rather, it only implies that the procedure would not perform its intended function due to its inadequate implementation by an experimental system that lacks sufficient system robustness.'' The procedure's definition can be in order, and the implementation can still fail. That sentence is the obstacle to reading a mathematical robustness theorem as a realization theorem.

A third notion comes from the philosophy of modeling, and it is different again. \citet{kuorikoski2007} describe derivational robustness as the stability of a modeling result under changes in modeling assumptions. They report Weisberg's robustness analysis as four steps: check whether a set of models imply a common result, look for a common structure, formulate a robust theorem, and ask which conditions defeat the link between the structure and the property. I have not re-read Weisberg, so those four steps are their report, not an independent exegesis. They also quote Orzack and Sober to the effect that ``the robustness of \(R\) is not by itself a reason to believe it,'' and they report Sugden's claim that robustness comparisons among models do not license the inductive leap from the models to the world. Kuorikoski, Lehtinen, and Marchionni think derivational robustness can still earn its keep inside a model, by showing that a result is not an artifact of a particular auxiliary assumption. They do not think it is, by itself, empirical confirmation. That limited defense is the most this paper will claim for the finite theorems, and only after saying which notion is in play.

The warning against transfer is then easy to apply, and it is the warning they attribute to Woodward. A proof that a procedure returns the same kind of answer for every input in a class is not a proof that several independent detections have converged, and it is not a proof that a causal relation in the world is stable. Calling all three ``robustness'' does not move the warrant from one to the next.

\section{What the finite theorems are robust about}
\label{sec:math}

A process function, in the sense of \citet{baumeler2016}, is a map \(w\) from a tuple of outputs to a tuple of inputs, on finite alphabets, such that every tuple of local functions has exactly one fixed point. The fixed point is the history for that tuple. The quantifier is the whole point. One history for one choice of local functions is not the condition.

Read as a procedure, the definition is a test. The input is a tuple of local functions. The intended function of the test is to return the history. Procedure robustness, in Karaca's sense, would be the capacity of that test to keep doing so under variation of the input. Universal consistency is the extreme case of that capacity inside the model: every input in the class of local functions is a success, and the success is uniqueness as well as existence. The identity-and-bit-flip wiring fails existence. The identity-and-copy wiring fails uniqueness. Both fail procedure robustness of this test. The three-bit rule \(w(a,b,c)=(\lnot b\land c,\ \lnot c\land a,\ \lnot a\land b)\) passes it, including for the three copying functions, whose only history is \((0,0,0)\). These are calculations about the test. They are not results of detections.

The canonical extension is a basis permutation defined for every \(w\), not only for process functions: \(U\lvert o,e\rangle=\lvert w(o)+e,o\rangle\). It is reversible as a relabeling whether or not the test succeeds. The further theorem says that the closed loop obtained by inserting local unitaries, on finite ancillas whose dimensions match, is unitary for every such insertion if and only if \(w\) is a process function. The forward direction uses an exclusivity lemma: distinct output tuples that a process function can produce must differ, at some party, in a way that makes a row of the loop orthogonal. The converse uses one family of local unitaries, the reversible encodings of the local functions, and observes that a column's squared norm equals the number of fixed points. The equivalence is therefore not the claim that some dilation or other is unitary. It is the claim that this family of procedures, now with unitary inputs rather than function inputs, preserves unitarity of the loop exactly when the classical test is procedure-robust in the sense above.

That is a derivational result in the narrow sense of \citet{kuorikoski2007}: a theorem inside a model, stable under a stated class of changes to the local maps. It is not Weisberg's four-step analysis, because there is no ensemble of independently constructed models here, and no robust theorem extracted from their common structure. There is one wiring, one test, and a universal quantifier. Stability under that quantifier is a strong form of insensitivity to the local maps. It is a weak form of derivational robustness if ``derivational robustness'' means agreement across a diverse set of modeling frameworks. The two should not be praised with one compliment.

The same section of the mathematics says what the theorem is not. Ancillas may begin entangled; the unitarity claim is about the loop operator, not about a product state. Instruments that are not unitary are covered only by a finite Stinespring dilation, padding, and a discard, and only because the unitary case is already settled \citep{stinespring1955}. Continuous alphabets are not covered. With infinitely many parties, exclusivity still gives at most one fixed point, and existence can fail. Whether every exclusive Baire-class-one self-map of an infinite product of bits has a fixed point is open, and unused. Finite explicit tables can be searched. The search shows that the condition is decidable for a given table. It is not a complexity theorem.

\section{Why the robustness is not evidence of a target}
\label{sec:evidence}

Suppose the test is procedure-robust. What may one conclude about a phenomenon outside the model?

\citet{kuorikoski2007} quote Orzack and Sober's claim that the robustness of a theorem \(R\) is not by itself a reason to believe \(R\), and they report Sugden's claim that comparisons among models do not license an inductive leap to the world. Their own reply is that robustness can show which assumptions a result does not depend on, and can thereby say which part of a model a later empirical test would be testing. That reply does not issue a free pass. It says that robustness earns a role in the design of a test. It does not replace the test.

Apply the reply here. The process-function theorem shows that unique solvability does not depend on which local functions are inserted, if it holds for all of them, and that unitarity of the canonical loop does not depend on which local unitaries are inserted, under the same condition. A later empirical test of a device could, in principle, be aimed at that stable feature rather than at a particular choice of gates. The theorem does not provide the test. It does not say which device, which records, or which interventions correspond to the local maps. Until those correspondences are stated, there is nothing for data to confirm or refute except the arithmetic.

Result robustness, as Karaca defines it, is not available as a substitute. There are not two independent means of detection whose outcomes agree. There is one definition and one proof. The three small wirings agree with the definition because they were computed from it. Agreement of that kind is not convergent validation. Karaca, following the line he attributes to Wimsatt and to Campbell and Fiske, treats independence of the means as essential to result robustness, and treats procedure robustness as necessary but not sufficient to rule out spurious convergence. A single derivation cannot fail that test by lacking independence of several derivations. It fails it by not being that kind of achievement at all.

There is a standing temptation to treat the universal quantifier as if it were an ensemble of models. Each local replacement is a different input, and the history exists in each case, so the existence of histories looks like a robust theorem in Weisberg's reported sense. The temptation should be resisted. In the report of Weisberg's steps, the models are different constructions, and the robust theorem concerns a property they share for a structural reason one then tries to isolate. Here the ``models'' are not different constructions. They are the input slots of one procedure. The common result is the procedure doing what it was defined to do. That is procedure robustness by stipulation when the definition is satisfied. It is not a discovery that several ways of modeling a target agree.

Hausman and Woodward's directed modularity is a different invariance again, and this paper does not relitigate it \citep{hausman1999}. Their condition, as used in a companion adaptation of the same mathematics, concerns equations that remain invariant when an intervention breaks an arrow. Process-function consistency concerns existence and uniqueness of a solution under replacement of local maps, including replacements that leave a cycle in place. A cyclic wiring can pass the second test and still not be modular in their directed sense. Nothing here claims they were wrong about directed modularity. The point of mentioning them is to refuse a transfer: invariance of solvability is not their invariance, and the warrant one might give their condition, as an interventionist constraint on causal order, does not become a warrant for implementation of a process function.

\section{Implementation}
\label{sec:impl}

Karaca's sentence about inadequate implementation is the right standard. A procedure is one thing. A system that carries it out is another. The system can fail to be robust, or fail to be the right system, while the procedure remains well defined.

The bare permutation \(U\) is the clearest case. It exists for every \(w\). It is a reversible map on a source-sink space. It does not, by existing, mean that a laboratory device performs the wiring step by step, and it does not mean that the wiring has a history. Bennett's logical reversibility, which is a different adaptation's topic, is a single-valued inverse of a machine's transition function. The permutation \(U\) is not that construction. Citing it as ``the system which implements the procedure'' would confuse a relabeling with a machine. Karaca's distinction predicts the confusion: one looks at a reversible mathematical object and calls it the implementation, because the procedure one cares about also has a reversibility-flavored theorem nearby.

The finite search is a second false implementation. Given an explicit table, one can enumerate output tuples and count fixed points. The search is itself a procedure, and for a finite table it is robust in a dull sense: it halts with the count. Its intended function is to evaluate the definition. The wiring's intended function, if one grants it an intended function at all, is to determine histories under local interventions. A procedure that checks a definition does not thereby realize the system the definition describes. The decidability of the check blocks a reading on which the histories are uncomputable. It does not install the wiring in a piece of hardware, and it does not yield a speedup for a problem the search was not given.

The unitary equivalence is a third false implementation, subtler because it quantifies over operations that look like physical gates. Unitarity under every local unitary is a fact about the loop operator in the model. To treat it as a realized computer one needs a correspondence between those unitaries and operations a device can undergo, and between the histories and records the device produces, and a reason to think the correspondence holds in the regime where the quantifier is being invoked. Karaca's remark says that failure of the system is not failure of the procedure. The converse, which he also leaves room for, is what matters here: success of the procedure is not success of the system. The theorem is entirely on the procedure's side.

A known numerical witness does not close the gap. A switchboard of three-qubit unitaries, fixed by a published table of phases, induces a conditional distribution \(p(i\mid o)\) in the computational basis. A linear functional \(L\) evaluates to \(1/4-(3/16)\sqrt{2}<0\) on that distribution. Nonnegativity of \(L\) on every three-bit process function is an external finite enumeration, not re-derived in this file. If the enumeration is granted, \(p\) is not a mixture of those process functions. The board is a variant of a known purification, not a new device. Read in another basis, a pure process of this family need not leave the hull. The calculation is a counterexample to a universality claim inside the model: purifiability by the canonical route, or by this board, does not restore membership in the convex hull, in the named basis. It is not a detection by independent means, and it is not a report that a laboratory ran the board.

Two-party classical processes are, in the literature, mixtures of causally ordered process functions. That classification is not re-proved here. It is cited only to say why the witness uses three parties \citep{oreshkov2012}. The number of parties is not a story about implementation either.

\section{Objections}
\label{sec:obj}

The first objection says that procedure robustness under every local replacement is so strong that the inductive leap is no longer a leap. If the result survives every intervention the model allows, nothing the model calls an intervention can be a confounding artifact, and the result is as secure as a theorem can be. The objection is right about artifacts inside the model and wrong about the leap. Orzack and Sober's point, as quoted by \citet{kuorikoski2007}, is that robustness is not by itself a reason to believe the robust claim as a claim about the world. Surviving every input in a class shows that the claim does not depend on which member of the class was chosen. It does not show that the class matches the interventions available in a target, or that the claim is true of the target. Strength of the quantifier increases security inside the model. It does not change the subject of the security.

The second objection says that unitarity is a physical property, so a unitary lift is already an implementation. Unitarity is a property of an operator in a Hilbert space. Whether a device's dynamics are given by that operator is a further claim. The bare permutation is unitary for wirings that have no history and for wirings that have two. If unitarity were implementation, those wirings would be implemented as well, including the ones that fail the solvability test. They are not, on any standard that uses the test as the condition for a process function. The objection proves too much, and what it proves is false to the theorem's own quantifiers.

The third objection says that the witness is empirical because it uses a quantum state and a Born-rule probability. The state and the probabilities are computed from a table of exponents and a formula. No sample was drawn. The negativity of \(L\) is an exact contraction of that table. Calling the contraction an experiment would make every evaluated formula an experiment. Karaca's result robustness requires means of detection. A formula is not a detector. The external enumeration of process functions is also not a detector. It is a finite check whose details are not in this file. A reader who withholds it should withhold the claim that \(p\) lies outside the hull, and should not replace the withheld check with an experimental story.

The fourth objection says that distinguishing procedure from system is a verbal delay. Once the procedure is specified completely, a system that meets the specification implements it, and the canonical unitary is the specification. The objection slides from a specification to one mathematical object that is related to the specification. The specification of a process function is the fixed-point condition. The canonical unitary meets a biconditional with that condition, under a stated class of tests. Meeting a biconditional is not being a physical system. Even if one constructed a device whose interventions and records were proved to match the condition, the proof of the match would be the implementation argument, and it is not among the theorems. Karaca's distinction is verbal only if ``procedure'' and ``system'' were already synonyms. His text uses them to mark a failure mode: the procedure can be fine and the implementation inadequate. That failure mode is available here, and the theorems do not close it.

The fifth objection says that this paper has invented a disagreement with Karaca by taking a definition from high-energy physics and applying it to a toy wiring. The application is an analogy, and it should be labeled as one. Karaca's intended function is the function of a procedure in an experiment, such as a data-selection cut. The solvability test is not a cut in ATLAS. What is being used is the invariance pattern: function maintained under input variation, distinguished from convergence of results, distinguished from the system that would run the procedure. If that pattern does not survive the move from his cases to this test, the argument fails. I think it survives, because nothing in the pattern depended on calorimeters. What does not survive is any claim about his ATLAS case itself. This paper does not re-interpret the data-selection procedure at the LHC.

\section{Conclusion}
\label{sec:end}

Universal consistency is a strong invariance. It says that a solvability test keeps its intended answer, one history, under every local replacement. The canonical unitary theorem says that a related test, unitarity of a closed loop, keeps its intended answer under every local unitary and every matched finite ancilla, and that the two invariances coincide for that canonical construction. Those are facts about procedures defined on a model. They are the sort of fact \citet{karaca2022} calls procedure robustness, transferred, with the limits in Section~\ref{sec:obj}, from experimental procedures to a mathematical test.

They are not the convergence of independent detections. They are not, by themselves, a reason to believe a claim about a target, in the sense of the warning \citet{kuorikoski2007} quote from Orzack and Sober and attribute, as a general caution, to Woodward. They are not an implementation. Karaca's distinction between a robust procedure and a system adequate to carry it out is the distinction the unitary permutation most often obscures. The permutation can be written down for a wiring the test fails. A search can decide the test without being the wiring. A negative value of a functional in one basis can separate a calculated distribution from a convex hull without being a measurement.

The paper does not offer a new robustness concept and does not correct Karaca's account of ATLAS. It uses his distinctions, and the caution in the modeling literature, to say what the existing finite theorems will support. They will support a claim of procedure-level invariance inside a stated model. They will not support a claim that a system has been realized, that a phenomenon has been explained by convergent detection, or that a reversibility-flavored operator is a machine. Those further claims need correspondences this mathematics does not state. Calling the invariance robustness does not supply the correspondences.

\appendix
\section{Account of generative assistance}
\label{app:ai}

The journal's author FAQ requires an AI statement, and, where the use is more than a simple translation, an account sufficient to reconstruct the steps: the tool, the date or date range, and the exact wording of the prompts. This appendix is the account. It is incomplete in ways that should not be papered over.

The tool was Grok Bot, a Cursor cloud agent. No model version identifier is available, and none is invented. The date of the drafting reported here is 8 November 2026. The operative instruction, in the user's words, was to write a manuscript for the \emph{British Journal for the Philosophy of Science} after the Physical Review A draft, as a sustained philosophy-of-science argument about operational robustness, scientific explanation, and implementation, not a theorem tour; to engage relevant primary literature; to avoid fabricated results; and to label the piece as an alternative adaptation if a new contribution could not be supported. Intermediate prompts in the same working session are not retained here as a complete verbatim log. A reader cannot reconstruct every step from this appendix alone. Human verification is pending. The agent is not an author.

The mathematical statements were taken from public drafts of the same finite results. No new proof was asked for or supplied. Funding, affiliation, and competing interests are not known and are not declared as absent. They are omitted from the review file because the journal requires acknowledgements to be deleted in full, not because they have been checked. A separate identity sheet, which is not part of the anonymized review text, records the author's name and the pending items. Because related drafts are already public, anonymity is not claimed. The title and abstract of this file are not the titles of those drafts. The public technical note is not cited by the author's name in this review file, which departs from the journal's instruction to cite one's own work in the third person under one's name. The departure is deliberate: the name would identify the author of the public note. The identity sheet records the departure.

\begin{thebibliography}{99}
\bibitem[Baumeler and Wolf(2016)]{baumeler2016} Baumeler, \"A., and S. Wolf. 2016. ``The Space of Logically Consistent Classical Processes without Causal Order.'' \textit{New Journal of Physics} 18:013036.
\bibitem[Hausman and Woodward(1999)]{hausman1999} Hausman, D. M., and J. Woodward. 1999. ``Independence, Invariance and the Causal Markov Condition.'' \textit{British Journal for the Philosophy of Science} 50:521--583.
\bibitem[Karaca(2022)]{karaca2022} Karaca, K. 2022. ``Two Senses of Experimental Robustness: Result Robustness and Procedure Robustness.'' \textit{British Journal for the Philosophy of Science} 73:279--298. \url{https://doi.org/10.1093/bjps/axy031}. Read from the text of the article as extracted for this draft.
\bibitem[Kuorikoski, Lehtinen, and Marchionni(2007)]{kuorikoski2007} Kuorikoski, J., A. Lehtinen, and C. Marchionni. 2007. ``Economics as Robustness Analysis.'' PhilSci-Archive preprint, 25 September 2007. \url{https://philsci-archive.pitt.edu/3550/}. This is the text that was read. A later journal version, if different, was not collated.
\bibitem[Oreshkov, Costa, and Brukner(2012)]{oreshkov2012} Oreshkov, O., F. Costa, and \v{C}. Brukner. 2012. ``Quantum Correlations with No Causal Order.'' \textit{Nature Communications} 3:1092.
\bibitem[Stinespring(1955)]{stinespring1955} Stinespring, W. F. 1955. ``Positive Functions on \(C^{*}\)-Algebras.'' \textit{Proceedings of the American Mathematical Society} 6:211--216.
\bibitem[Woodward(2006)]{woodward2006} Woodward, J. 2006. ``Some Varieties of Robustness.'' \textit{Journal of Economic Methodology} 13:219--240. Not re-read. Cited only as reported by Kuorikoski, Lehtinen, and Marchionni.
\end{thebibliography}
\end{document}
